\section{Clustered and Calibrated Inference}
\label{app:clustered-inference}

This appendix isolates the part of the theory stack where clustered standard
errors, calibration, and the final TreePO bound object meet. The earlier
appendices explain how to estimate local violation rates and how to obtain
confidence radii. The purpose here is narrower: once a clustered confidence
interval has been constructed, how does it feed into the bound we actually
report?

\subsection{A Generic Interval Lemma}

\begin{lemma}[Confidence interval from error control]
\label{lem:clustered-ci-membership}
Let $\theta$ be a scalar target, let $\widehat{\theta}$ be an estimator, let
$\widehat{\mathrm{SE}}\ge 0$ be a standard error estimate, and let $z\ge 0$.
If
\[
|\theta-\widehat{\theta}| \le z\,\widehat{\mathrm{SE}},
\]
then
\[
\theta \in
\left[
\widehat{\theta}-z\,\widehat{\mathrm{SE}},
\widehat{\theta}+z\,\widehat{\mathrm{SE}}
\right].
\]
\end{lemma}

\begin{proof}
The displayed absolute-value bound is equivalent to the pair of inequalities
\[
-z\,\widehat{\mathrm{SE}}
\le
\theta-\widehat{\theta}
\le
z\,\widehat{\mathrm{SE}}.
\]
Adding $\widehat{\theta}$ throughout yields
\[
\widehat{\theta}-z\,\widehat{\mathrm{SE}}
\le
\theta
\le
\widehat{\theta}+z\,\widehat{\mathrm{SE}},
\]
which is exactly the claimed interval membership.
\end{proof}

\paragraph{Event form.}
If an event $E$ implies
$|\theta-\widehat{\theta}| \le z\,\widehat{\mathrm{SE}}$,
then the same event implies coverage of the corresponding confidence interval.
This is the only input needed to turn any external concentration theorem or
high-probability certification event into a confidence-interval statement.

\subsection{Judge-Bias Envelope}

Let $b_\star$ denote the true judge bias, let $\widehat b$ be the estimated
bias, and let $\widehat{\mathrm{SE}}_b$ be either the usual or the clustered
standard error for that bias estimate.

\begin{corollary}[Absolute-bias envelope from confidence-interval membership]
\label{cor:abs-bias-envelope}
If
\[
b_\star \in
\left[
\widehat b-z\,\widehat{\mathrm{SE}}_b,\,
\widehat b+z\,\widehat{\mathrm{SE}}_b
\right],
\]
then
\[
|b_\star|
\le
|\widehat b| + z\,\widehat{\mathrm{SE}}_b.
\]
\end{corollary}

\begin{proof}
By Lemma~\ref{lem:clustered-ci-membership},
\[
|b_\star-\widehat b| \le z\,\widehat{\mathrm{SE}}_b.
\]
Now apply the triangle inequality:
\[
|b_\star|
\le
|b_\star-\widehat b| + |\widehat b|
\le
z\,\widehat{\mathrm{SE}}_b + |\widehat b|.
\]
This is the claimed envelope.
\end{proof}

\paragraph{Why this matters.}
The calibration layer uses exactly this quantity. The reported absolute-bias
radius is not an additional assumption once the confidence interval has been
constructed; it is a deterministic corollary of interval membership. The same
argument applies with the clustered bias standard error, so document-level
dependence simply changes which $\widehat{\mathrm{SE}}_b$ enters the formula.

\subsection{\texttt{computeDSLBound} from a Joint Event}

Write the judge-side TreePO estimate as $\widehat G_{\mathrm{tree}}$, its
clustered standard error as $\widehat{\mathrm{SE}}_{\mathrm{tree}}$, and the
calibration radius as $B_{\mathrm{cal}}$. The upper endpoint reported by the
system is
\[
U
:=
\widehat G_{\mathrm{tree}}
 + z\,\widehat{\mathrm{SE}}_{\mathrm{tree}}
 + B_{\mathrm{cal}}.
\]
This is exactly the quantity returned by \texttt{computeDSLBound} when the
estimate is nonnegative, as it is for the union-bound based tree estimate.

\begin{theorem}[Joint interval-plus-calibration validity]
\label{thm:joint-interval-calibration}
Assume
\[
G_{\mathrm{judge}}
\in
\left[
\widehat G_{\mathrm{tree}} - z\,\widehat{\mathrm{SE}}_{\mathrm{tree}},\,
\widehat G_{\mathrm{tree}} + z\,\widehat{\mathrm{SE}}_{\mathrm{tree}}
\right]
\]
and
\[
|G_{\mathrm{oracle}} - G_{\mathrm{judge}}| \le B_{\mathrm{cal}}.
\]
If $\widehat G_{\mathrm{tree}} \ge 0$, then
\[
|G_{\mathrm{oracle}}| \le U.
\]
\end{theorem}

\begin{proof}
The interval membership assumption gives
\[
G_{\mathrm{judge}}
\le
\widehat G_{\mathrm{tree}} + z\,\widehat{\mathrm{SE}}_{\mathrm{tree}}.
\]
Because the interval is centered at a nonnegative quantity, the same upper
endpoint also bounds the absolute value:
\[
|G_{\mathrm{judge}}|
\le
\widehat G_{\mathrm{tree}} + z\,\widehat{\mathrm{SE}}_{\mathrm{tree}}.
\]
Now apply the triangle inequality:
\[
|G_{\mathrm{oracle}}|
\le
|G_{\mathrm{oracle}}-G_{\mathrm{judge}}| + |G_{\mathrm{judge}}|
\le
B_{\mathrm{cal}}
 + \widehat G_{\mathrm{tree}}
 + z\,\widehat{\mathrm{SE}}_{\mathrm{tree}}
= U.
\]
\end{proof}

\begin{theorem}[Joint interval-plus-calibration validity with oracle error]
\label{thm:joint-interval-calibration-oracle}
Assume the hypotheses of
Theorem~\ref{thm:joint-interval-calibration}, and in addition suppose
\[
|G_{\mathrm{true}} - G_{\mathrm{oracle}}| \le B_{\mathrm{oracle}}.
\]
Then
\[
|G_{\mathrm{true}}| \le U + B_{\mathrm{oracle}}.
\]
\end{theorem}

\begin{proof}
By Theorem~\ref{thm:joint-interval-calibration},
$|G_{\mathrm{oracle}}| \le U$. One more triangle inequality gives
\[
|G_{\mathrm{true}}|
\le
|G_{\mathrm{true}}-G_{\mathrm{oracle}}| + |G_{\mathrm{oracle}}|
\le
B_{\mathrm{oracle}} + U.
\]
\end{proof}

\paragraph{Interpretation.}
The theorem has a useful modular form. A clustered confidence interval controls
the judge-side tree estimate. Calibration controls judge-vs-oracle error.
Optional oracle measurement control handles oracle-vs-truth error. The final
reported upper bound is obtained by adding exactly those radii, with no further
statistical trick.
